%!TEX program = uplatex
\RequirePackage{plautopatch}
\documentclass[uplatex,dvipdfmx]{jlreq}
\usepackage{amsmath,amssymb}

\pagestyle{empty}
\begin{document}

\begin{center}
  {\large 概要}
\end{center}

\vspace{20pt}

\noindent
{\large\bfseries 応用特異点論概説}
\hfill
{\large\bfseries 泉屋　周一（北大・理）}

\medskip

今年（1998年）の夏に北海道大学理学部数学科のパイロット事業「高校生のための夏期講座」において、
「３次方程式とコーヒーカップの底」と題した講演を３時間ほど行った。
そこでは、３次方程式の標準形における判別曲線がカスプ型となること、
および乳白色のコーヒーカップの底にあたる光から出来る焦線の形が同じくカスプ型になることが観察され、
その両者にはいかなる関係があるかということを高校一年生程度の数学知識を仮定して解説した。

当講演では、同様な話題から出発して、
日常生活で観察される様々なカスプ型について特異点論的にはどのように解釈されるのかを解説したい。
また、簡単に可微分写像の特異点論の歴史を概観し、石川氏による基礎理論の解説につなげたい。
これにより、基礎理論の必要性と応用の必然性の理解の補助となるであろう。

当日、レーザーポインターを使って、実際にカスプ型をスクリーン上に映し出す実験も試みたい。

\end{document}